\section{Discussion}\label{sec:discussion}

\begin{guide}
\textbf{About 1.5 pages.} This is where marks for understanding are earned. Explain \emph{why} the results came out this way, and refer to the tables rather than repeating them. Write four paragraphs:
\begin{enumerate}
  \item \textbf{Risk, not prediction.} Direction is a coin flip, yet drawdown fell in both periods. The mechanism: the engine controls \emph{when} and \emph{how much} to be exposed, through fewer trades, smaller and better-placed positions on volatile days, and staying out on high jump-risk days. Volatility clusters and is forecastable; daily returns are close to a random walk.
  \item \textbf{H1, H2 and the volatility target.} Extra features and attention barely help because next-day returns carry very little signal, the dataset is small (about 1{,}500 training days), and deeper RNNs train unstably. The 14-day target is a lesson for others: always check a forecast against persistence.
  \item \textbf{The liquidation result.} Volatility-widened stops sat beyond the 5\% liquidation line at 20x. The fix is to cap the stop inside $1/L$. It was deliberately not applied, because changing the engine after seeing test results would be tuning on the test set.
  \item \textbf{Deviations and limitations.}
  \begin{itemize}
    \item Deviations from the proposal: the single-day volatility target, the added 2026 forward test, and the extra ablations and bootstrap tests.
    \item Limitations: one asset over a short test period with high seed-to-seed variance; no slippage or CFD financing costs \tbc{check}; balances in rand ignore USD/ZAR moves; the volatility forecast is biased high; absolute returns are small (+4.5\% in 2025, +2.4\% in 2026).
  \end{itemize}
\end{enumerate}
\end{guide}

\tbc{}
